\section{Matched quality evaluation for image
attacks}\label{record-docs-matched-image-quality-md-matched-quality-evaluation-for-image-attacks}

This extends all preserved model-weight attacks with the paired metrics
used for NFPA, and adds the earlier class-text CLIP metric to every NFPA
arm. It reads existing PNGs; no model training, image generation or
attack is repeated.

For each of RAR-XL and UViT-H/DC-AE, all 512 original watermarked
reference PNGs were verified byte-for-byte identical across the earlier
weight-attack and NFPA runs. Weight outputs are paired by original
generation index and checked against the same ImageNet class; NFPA pairs
are recovered using the hash of its original input PNG. Missing,
duplicate or mismatched pairs cause an error. Native resolutions are
retained for paired pixel/perceptual metrics.

Every condition now receives the same metrics:

\begin{itemize}
\tightlist
\item
  LPIPS with AlexNet v0.1, using input pixels scaled to {[}-1, 1{]}.
  Lower is better for preservation of the reference image.
\item
  SSIM at native resolution with unit pixel range. Higher is better.
\item
  PSNR with unit pixel range, capped at 120 dB for identical images,
  matching the existing NFPA evaluator.
\item
  CLIP image-image cosine between original and candidate embeddings.
\item
  CLIP image-class-text cosine against
  \texttt{\seqsplit{a photo of a <ImageNet class>}} using torchvision's
  canonical 1,000-class ordering. This reproduces the earlier alignment
  metric and adds it to NFPA.
\end{itemize}

The pinned CLIP ViT-B/32 checkpoint and LPIPS weights/cache are shared
with the existing evaluation. All outputs, data tables, relevant model
assets and source are hashed. Previous class-alignment scores and NFPA
paired scores are checked against recomputation (maximum permitted
absolute discrepancy 0.001); actual maximum differences are saved.
Results include per-image scores, means, percentiles and 95\% bootstrap
intervals over image pairs. Detection rates are copied from the original
frozen owner evaluation without recalibration.

Nine weight conditions include the unmodified supplied model, all four
merge settings (including clean replacement), and both 500/2,500-step
LoRA/full updates. Four output conditions include resize-only and NFPA
delta 0/20/40. This gives 13 conditions with 512 pairs per model, 13,312
candidate/reference pairs overall.

Paired distances measure how much a generated sample changes. A model
whose weights changed can generate a different valid image for the same
class/seed. Consequently, high LPIPS or low SSIM alone does not
establish worse generation quality. Class alignment is a complementary
semantic proxy; neither metric certifies exact identity, factual content
or overall visual quality.

\subsection{Commands}\label{record-docs-matched-image-quality-md-commands}

\begin{Verbatim}
PYTHONPATH=. uv run --no-sync python scripts/compare_image_attack_quality.py prepare
PYTHONPATH=. uv run --no-sync python scripts/compare_image_attack_quality.py submit RUN/manifest.json
PYTHONPATH=. uv run --no-sync python -m pytest ai_content_detector/tests/test_image_quality.py -q
\end{Verbatim}

Dependencies already belong to
\texttt{\seqsplit{requirements\_watermarking.in}} and the existing
research environment. One restartable H100 dev-QoS job runs per image
family. Metrics are saved every eight pairs; jobs checkpoint and requeue
on the Slurm advance signal. Source image and prior result files are not
overwritten.

Prepared run: \texttt{\seqsplit{20261104T153348508006Z}}. Submitted
one-GPU jobs \texttt{\seqsplit{1958522}} (RAR) and
\texttt{\seqsplit{1958524}} (DC-AE). Four targeted CPU tests passed,
including pairing, PNG integrity/native geometry, PSNR range and
corresponding-label CLIP lookup. GPU completion and numerical agreement
are recorded in per-family receipts.

Both full jobs completed successfully. See the
\href{records/docs/matched_image_quality_results.tex}{combined matched
results}. Every condition has 512 pairs; maximum prior-metric
discrepancy was 2.3842e-7 on each GPU. The real-metric CPU smoke passed
on both targets as well.
